% !TeX root = ../main.tex
\chapter{Simulation-Based Validation}
\label{ch:simulation}

The hardware loop of \cref{ch:integration} mixes three contributions to
the recorded \ee{} error: the residual of the optimiser on its own
kinematic model, the lag and clamping of the MRT bridge, and the
odometry-plus-forward-kinematics estimate of where the robot is. A
kinematic closed loop that replaces the bridge with an ideal rollout
removes the last two, so a fault in the reference, the time axis or the
cost function can be seen before it is blamed on the plant. That loop is
the dummy MRT of \ocs{}. \Cref{sec:sim-methodology} states what it can
show and defines the metrics used here and in \cref{ch:hardware}.
\Cref{sec:sim-dummy} describes the closed \ee{} path that the dummy
exercises and the checks that a session must pass before the same files
are taken to the robot.

% ---------------------------------------------------------------------
\section{Validation Methodology and Evaluation Metrics}
\label{sec:sim-methodology}

The dummy MRT is not a second controller. It is the production \mpc{}
node and the production target publisher, closed through
\texttt{MRT\_ROS\_Dummy\_Loop} instead of the hardware bridge
(\cref{sec:hw-interfaces}). The dummy integrates the latest policy with
the same kinematic rollout the solver uses, at the configured MRT rate of
\qty{100}{\hz}, and feeds the resulting state back as the next
observation. There is no policy-evaluation lag, no actuator dynamics, and
no wheel-odometry drift. The \ee{} error that remains is the optimiser's
residual against its own model: input costs and velocity limits prevent a
perfect match to the interpolated reference, and a finite horizon cannot
see the whole path.

That residual is useful. If the dummy already leaves the commanded path,
the later hardware error is not an execution problem. If the dummy
tracks and the hardware does not, the gap is in the bridge, the
actuators, or the state estimate, which is the comparison drawn in
\cref{ch:hardware}. The dummy does not speak to real-time behaviour on
the IPC: solve time, policy exhaustion and the observation-age guard
never appear, because the observation is produced in the same process
tree as the solver.

A second simulated-hardware mode exists in the \hero{} stack: mock
\texttt{ros2\_control} interfaces that echo commands into states while
still running the MRT bridge~\cite{ros2control}. That mode keeps the
timing and clamping of \cref{sec:hw-interfaces} and is used as a
stepping stone onto the physical platform. It is not the subject of this
chapter.

\subsection{Internal metrology}
\label{sec:sim-metrology}

Desired and attained \ee{} poses are both computed inside the waypoint
node from the same \mpc{} observation. The desired pose is the linear
interpolation of the commanded knots at \texttt{observation.time}, which
is the interpolation \ocs{} itself uses for the position part of the
reference. The attained pose is Pinocchio forward kinematics of that
observation, with \texttt{tool0} as the \ee{} frame
(\cref{sec:hero-platform}). The error signal therefore carries no
cross-topic time-alignment residual. The node publishes it at
\qty{20}{\hz} on \texttt{/ee\_waypoint\_target/ee\_metrics}.

On the dummy this is a consistent kinematic measurement. On hardware the
same pipeline reconstructs the \ee{} from wheel odometry and joint
encoders. Drift of the base pose then contaminates every later metric.
The figures and tables of this chapter, and the internal traces of
\cref{ch:hardware}, are therefore internal metrology. An ISO-conforming
accuracy statement would require an external pose of the \ee{}, for
example from an OptiTrack system, with the uncertainty budget of
ISO~9283:1998 cl.~6.5~\cite{ISO9283}. That measurement is not part of
the present experiments.

\subsection{Time-aligned pose error}
\label{sec:sim-pose-error}

Let $\symbf{p}_\mathrm{d}(t)$ and $\symbf{p}_\mathrm{a}(t)$ be the
desired and attained positions in the \texttt{world} frame, and
$q_\mathrm{d}(t)$, $q_\mathrm{a}(t)$ the corresponding unit quaternions.
The position and orientation errors used throughout are
\begin{equation}
  e_{\mathrm{p}}(t) = \lVert \symbf{p}_\mathrm{d}(t) -
    \symbf{p}_\mathrm{a}(t) \rVert,
  \qquad
  e_{\mathrm{o}}(t) = 2\,\mathrm{atan2}\bigl(\lVert \symbf{v} \rVert,
    \lvert w \rvert\bigr),
  \label{eq:ee-errors}
\end{equation}
where $(w,\symbf{v})$ are the scalar and vector parts of
$q_\mathrm{d}\,q_\mathrm{a}^{-1}$ and $e_{\mathrm{o}}$ is folded into
$[0,\pi]$ so that the quaternion double cover does not appear as a
$2\pi$ jump. On a measured edge of the path the position error is split
into an along-track component $e_{\parallel} = \symbf{e}\cdot\symbf{\tau}$
and a cross-track remainder $e_{\perp} = \lVert \symbf{e} -
e_{\parallel}\,\symbf{\tau} \rVert$, with $\symbf{\tau}$ the unit
tangent of that edge. The ratio
\begin{equation}
  t_{\mathrm{lag}} = e_{\parallel} / v_{\mathrm{ref}}
  \label{eq:lag-eff}
\end{equation}
converts the along-track error into an equivalent time, with
$v_{\mathrm{ref}}$ the commanded path speed. On hardware,
$t_{\mathrm{lag}}$ is the quantity that can be compared with the
policy-lag diagnostic of the MRT bridge
(\cref{sec:hw-challenges}); the dummy has no such lag, so a large
$t_{\mathrm{lag}}$ there would mean the optimiser is chronically behind
its own reference.

These samples are only admitted on the measured path. The waypoint node
tags every metric with a segment identifier: \texttt{segment\_id} $\ge
1$ is an edge of the closed box; $0$ is the approach from the current
\ee{} to the entry point; $-10$, $-11$ and $-12$ are the lead-in,
lead-out and optional park legs of ISO~9283 cl.~8.2~\cite{ISO9283}.
Post-processing keeps $\texttt{segment\_id} \ge 1$. The transients of
starting, stopping and parking do not enter $e_{\mathrm{p}}$,
$t_{\mathrm{lag}}$ or the path metrics below.

\subsection{Path metrics: Hausdorff distance and ISO path accuracy}
\label{sec:sim-path-metrics}

Time-aligned error assumes that the commanded and attained poses share a
clock and a speed profile. That is true for the \ocs{} loop, where the
reference is a time-parameterised \texttt{TargetTrajectories} message.
It is not true of a decoupled Nav2 plus MoveIt baseline, which would
reach the same corners with a different velocity profile and with
pauses while the base repositions. A later comparison of the two
pipelines therefore needs a metric that treats the command as a
geometric path, not as a trajectory. The (symmetric) Hausdorff distance
between the commanded point set $A$ and the attained point set $B$ is
\begin{equation}
  h(A,B) = \max_{a \in A} \min_{b \in B} \lVert a-b \rVert,
  \qquad
  H(A,B) = \max\bigl(h(A,B),\, h(B,A)\bigr).
  \label{eq:hausdorff}
\end{equation}
The directed distance $h(B,A)$ is small whenever every attained sample
lies near the command, including the degenerate case of a robot that
stops halfway; $h(A,B)$ penalises shortcuts and unfinished laps. Both
directions are required. The classical $\max$ is sensitive to a single
spike, so results are reported as the symmetric mean
\begin{equation}
  \overline{H}(A,B)
    = \frac{1}{2|A|}\sum_{a\in A}\min_{b\in B}\lVert a-b \rVert
    + \frac{1}{2|B|}\sum_{b\in B}\min_{a\in A}\lVert a-b \rVert
  \label{eq:mean-hausdorff}
\end{equation}
together with the 95th percentile of the nearest-neighbour distances and
the raw $H(A,B)$. The command path is resampled at a spacing well below
the expected error (one millimetre is enough for the box of
\cref{sec:sim-dummy}) before the nearest-neighbour search, otherwise
four or eight corners would inflate every distance.

ISO~9283 defines path accuracy $\mathrm{AT}_{p}$ along a commanded path
as the maximum, over the path, of the distance between that command and
the barycentre of $n$ attained cycles~\cite{ISO9283}. With $n=1$,
$\mathrm{AT}_{p}$ reduces to the maximum deviation of a single attained
path from the command, which is close in spirit to $h(B,A)$ but still
parameterised by the command, not by an unordered cloud. Pose accuracy
$\mathrm{AP}_{p}$ (cl.~7.2.1) is the distance from each commanded corner
to the barycentre of the attained poses at that corner; it is the
fair common metric once a decoupled baseline exists, because that
baseline has no continuous rectangular path to score with
$\mathrm{AT}_{p}$. Cycle counts and the three mandatory speed levels of
ISO Table~3 are not exercised in this chapter. One slow lap on the dummy
is a pipeline check, not an ISO test.

Task completion time, the wall-clock duration of the measured lap, is
recorded alongside the geometric scores. Hausdorff distance ignores
speed; a slow but accurate motion would otherwise look like a win.

\Cref{tab:sim-metrics} summarises which quantity answers which question.
Dummy sessions use the first two rows as live diagnostics in RViz and in
the metric topic. Hausdorff and $\mathrm{AT}_{p}$ are computed
offline from the same bag export that \cref{ch:hardware} uses.

\begin{table}[htbp]
  \centering
  \caption[Metrics used in the closed-path experiments]{%
    Metrics used in the closed-path experiments. Dummy sessions produce
    the first two rows as a planner residual. Hausdorff distance and
    $\mathrm{AT}_{p}$ are evaluated offline on samples with
    $\texttt{segment\_id} \ge 1$.}
  \label{tab:sim-metrics}
  \begin{tabularx}{\linewidth}{@{}lXl@{}}
    \toprule
    Quantity & Question & Notes \\
    \midrule
    $e_{\mathrm{p}}$, $e_{\mathrm{o}}$
      & How far is the \ee{} from the command at this instant?
      & \cref{eq:ee-errors}; same observation for both poses \\
    $e_{\parallel}$, $e_{\perp}$, $t_{\mathrm{lag}}$
      & Is the error along-track lag or cross-track deviation?
      & measured edges only; \cref{eq:lag-eff} \\
    $H$, $\overline{H}$, $P_{95}$
      & Do the two paths have the same shape?
      & time-free; both directions required \\
    $\mathrm{AT}_{p}$
      & ISO path accuracy of the whole-body loop
      & cl.~8.2; whole-body pipeline only \\
    $\mathrm{AP}_{p}$
      & Did the \ee{} reach the corners?
      & cl.~7.2.1; common with a pose-to-pose baseline \\
    $T$
      & How long did the measured lap take?
      & wall clock; Hausdorff does not capture this \\
    \bottomrule
  \end{tabularx}
\end{table}

% ---------------------------------------------------------------------
\section{Kinematic Closed-Loop Simulation (Dummy MRT)}
\label{sec:sim-dummy}

The dummy launch starts, on one machine, the URDF exporter, the \mpc{}
node, the dummy MRT, the robot-state publisher, RViz, and the waypoint
target node. The interactive-marker target used in the \ocs{} demo is
switched off, so the waypoint node is the only publisher on the \mpc{}
target topic. The task file is \texttt{task\_experiment\_a.info}, the
same file the hardware host launch loads, and the waypoint file is
\texttt{experiment\_a\_waypoints.yaml}. What is being checked is the
configuration that will run on the robot, not a separate simulation
tune.

\Cref{fig:sim-dummy-stack} shows the resulting graph. Compared with
\cref{fig:architecture} there is no Zenoh boundary and no
\texttt{ros2\_control} layer. The dummy loop calls
\texttt{rolloutPolicy} on the kinematic model, advances
\texttt{observation.time} by $\qty{10}{\milli\second}$, and publishes
the new observation. The waypoint node consumes that observation both
to build knot times and to evaluate \cref{eq:ee-errors}.

\begin{figure}[htbp]
  \centering
  \begin{tikzpicture}[
      font=\footnotesize,
      box/.style={draw, rounded corners=1pt, align=center,
                  inner sep=4pt, minimum height=9mm},
      machine/.style={draw, dashed, rounded corners=2pt, inner sep=7pt},
      flow/.style={-{Latex[length=2mm]}, thick},
      lbl/.style={font=\scriptsize, align=center},
    ]
    \node[box, minimum width=34mm] (wp) at (0,0)
      {\texttt{ee\_waypoint\_target}\\reference + metrics};
    \node[box, minimum width=30mm] (mpc) at (5.2,0)
      {\mpc{} node\\(SLQ solver)};
    \node[box, minimum width=36mm] (dummy) at (5.2,-3.2)
      {dummy MRT\\kinematic rollout};
    \node[box, minimum width=28mm] (rviz) at (0,-3.2)
      {RViz\\commanded vs.\ attained};
    \node[machine, fit=(wp)(mpc)(dummy)(rviz),
          label={[lbl]above:single machine}] {};
    \draw[flow] (wp) -- (mpc)
      node[lbl, midway, above] {\texttt{\_mpc\_target}};
    \draw[flow] ([xshift=-4mm]dummy.north) -- ([xshift=-4mm]mpc.south)
      node[lbl, pos=0.45, left=1mm] {observation\\\qty{100}{\hz}};
    \draw[flow] ([xshift=4mm]mpc.south) -- ([xshift=4mm]dummy.north)
      node[lbl, pos=0.45, right=1mm] {policy};
    \draw[flow] (dummy.west) -- (rviz.east);
    \draw[flow] (wp.south) -- (rviz.north)
      node[lbl, midway, left=1mm] {markers,\\trace, $e_{\mathrm{p}}$};
  \end{tikzpicture}
  \caption[Dummy-MRT closed loop for Experiment~A]{%
    Dummy-MRT closed loop used to accept the Experiment~A configuration.
    The hardware bridge of \cref{fig:architecture} is replaced by a
    kinematic rollout of the latest policy; the waypoint node is
    otherwise unchanged.}
  \label{fig:sim-dummy-stack}
\end{figure}

\subsection{Commanded path}
\label{sec:sim-path}

The commanded path is an eight-corner rectangular parallelepiped whose
edges are aligned with the world axes after the box is landed
(\cref{fig:sim-box,tab:sim-waypoints}). ISO~9283 cl.~6.8.1 replaces a
cube by a parallelepiped when the travel along one axis is much smaller
than along the others~\cite{ISO9283}; here $z$ is an arm coordinate and
$x$, $y$ can recruit the base. The implemented box is smaller than the
workspace-filling $2.0\,\mathrm{m}\times 2.0\,\mathrm{m}\times
0.40\,\mathrm{m}$ parallelepiped discussed as a design target. Its
horizontal extents are \qty{0.60}{\meter} by \qty{0.80}{\meter} and its
height is \qty{0.25}{\meter}, with the lower face at
\qty{1.30}{\meter} and the upper face at \qty{1.55}{\meter}. The
perimeter is \qty{4.50}{\meter}.

Horizontal offsets are stored relative to the \ee{} pose at trigger
time (\texttt{origin\_frame: ee}); the $z$ values are absolute world
heights. Anchoring on the current \texttt{tool0} keeps the first corner
next to the standby pose, so the approach gate is a short motion rather
than a reach across the cell; the first offset is only
\qty{0.41}{\meter}. The far corner lies \qty{1.35}{\meter} from
that same \ee{}, which is enough to pull the base on the long edges
without demanding a $2\,\mathrm{m}$ translation on every lap. The
standby joint vector written in \texttt{task\_experiment\_a.info} (dummy
only; the hardware bridge ignores it) is the MoveIt named state
\texttt{bent}: shoulder pan, lift, elbow and wrist~1 at \ang{-90},
wrist~2 at \ang{+90}, wrist~3 at zero.

The visit order is a Hamiltonian cycle on the eight corners in which
every consecutive pair differs in a single coordinate:
\begin{equation*}
  \mathrm{P1}
    \xrightarrow{x}\mathrm{P2}
    \xrightarrow{y}\mathrm{P3}
    \xrightarrow{x}\mathrm{P4}
    \xrightarrow{z}\mathrm{P5}
    \xrightarrow{x}\mathrm{P6}
    \xrightarrow{y}\mathrm{P7}
    \xrightarrow{x}\mathrm{P8}
    \xrightarrow{z}\mathrm{P1}.
\end{equation*}
Six edges are horizontal and two are vertical. The vertical pair is a
pure arm motion in $z$; the base cannot contribute to it. Closing the
loop on a downward $z$ edge has a consequence for the ISO lead-out,
taken up below.

All eight knots share one orientation,
$\mathrm{RPY}=(\pi,\,0,\,-\pi/2)$, which is the \texttt{tool0} pose of
\texttt{bent} given the \ang{-90} mount of the arm
(\cref{sec:hero-platform}). A constant quaternion is required because
\texttt{TargetTrajectories::getDesiredState} interpolates quaternion
coefficients component-wise, not by spherical interpolation. A time-varying
reference orientation would inject interpolation error into
$e_{\mathrm{o}}$ that is not the controller's.

\begin{figure}[htbp]
  \centering
  \begin{tikzpicture}[
      font=\footnotesize,
      >=Latex,
      boxedge/.style={thick},
      zedge/.style={thick, dashed},
      lead/.style={thick, RWTHOrange100},
      pnt/.style={circle, fill=RWTHBlau100, inner sep=1.2pt},
      lbl/.style={font=\scriptsize},
    ]
    % scale: 1 m = 4.2 cm; x right = robot forward, y up = robot left
    \def\s{4.2}
    % bottom face z = 1.30
    \coordinate (P1) at ({0.40*\s}, {-0.10*\s});
    \coordinate (P2) at ({1.00*\s}, {-0.10*\s});
    \coordinate (P3) at ({1.00*\s}, {-0.90*\s});
    \coordinate (P4) at ({0.40*\s}, {-0.90*\s});
    % top face, offset to suggest height
    \coordinate (off) at (0.55, 0.70);
    \coordinate (P8) at ($(P1)+(off)$);
    \coordinate (P7) at ($(P2)+(off)$);
    \coordinate (P6) at ($(P3)+(off)$);
    \coordinate (P5) at ($(P4)+(off)$);
    \coordinate (Lin) at ({0.25*\s}, {-0.10*\s});
    % top rectangle (dashed) then bottom (solid), then verticals
    \draw[boxedge, dashed] (P8) -- (P7) -- (P6) -- (P5) -- cycle;
    \draw[boxedge] (P1) -- (P2) -- (P3) -- (P4) -- cycle;
    \draw[zedge] (P1) -- (P8);
    \draw[zedge] (P2) -- (P7);
    \draw[zedge] (P3) -- (P6);
    \draw[zedge] (P4) -- (P5);
    % measured Hamiltonian: P1-P2-P3-P4-P5-P6-P7-P8-P1
    \draw[->, very thick, RWTHBlau100] (P1) -- (P2);
    \draw[->, very thick, RWTHBlau100] (P2) -- (P3);
    \draw[->, very thick, RWTHBlau100] (P3) -- (P4);
    \draw[->, very thick, RWTHBlau100] (P4) -- (P5);
    \draw[->, very thick, RWTHBlau100] (P5) -- (P6);
    \draw[->, very thick, RWTHBlau100] (P6) -- (P7);
    \draw[->, very thick, RWTHBlau100] (P7) -- (P8);
    \draw[->, very thick, RWTHBlau100] (P8) -- (P1);
    \draw[lead, ->] (Lin) -- (P1);
    \foreach \p/\n/\a in {
      P1/P1/above, P2/P2/above, P3/P3/below,
      P4/P4/below, P5/P5/left, P6/P6/right,
      P7/P7/right, P8/P8/left} {
      \node[pnt] at (\p) {};
      \node[lbl, \a] at (\p) {\n};
    }
    \node[pnt, fill=RWTHOrange100] at (Lin) {};
    \node[lbl, above left] at (Lin) {lead-in};
    \node[lbl, anchor=west] at (5.0, -3.6)
      {solid: $z=\qty{1.30}{\meter}$; dashed: $z=\qty{1.55}{\meter}$};
    \node[lbl, anchor=west] at (5.0, -4.05)
      {$x$ forward (right), $y$ left (up)};
  \end{tikzpicture}
  \caption[Eight-corner commanded path for Experiment~A]{%
    Commanded parallelepiped in the \ee{}-anchored frame of
    \texttt{experiment\_a\_waypoints.yaml}. Blue arrows are the
    measured Hamiltonian cycle (segment identifiers $1$--$8$). The
    orange stub is the ISO~8.2 lead-in; the lead-out continues
    \qty{0.15}{\meter} down from P1 along the closing $z$ edge and is
    omitted here.}
  \label{fig:sim-box}
\end{figure}

\begin{table}[htbp]
  \centering
  \caption[Waypoint offsets of the Experiment~A box]{%
    Waypoint offsets of the Experiment~A box, copied from
    \texttt{experiment\_a\_waypoints.yaml}. $\Delta x$, $\Delta y$ are
    applied in the base yaw at trigger time about the current \ee{};
    $z$ is an absolute world height.}
  \label{tab:sim-waypoints}
  \begin{tabularx}{\linewidth}{@{}lXXXX@{}}
    \toprule
    Corner & $\Delta x$ / \si{\meter} & $\Delta y$ / \si{\meter}
      & $z$ / \si{\meter} & To next \\
    \midrule
    P1 & 0.40 & $-0.10$ & 1.30 & $+x$, \qty{0.60}{\meter} \\
    P2 & 1.00 & $-0.10$ & 1.30 & $-y$, \qty{0.80}{\meter} \\
    P3 & 1.00 & $-0.90$ & 1.30 & $-x$, \qty{0.60}{\meter} \\
    P4 & 0.40 & $-0.90$ & 1.30 & $+z$, \qty{0.25}{\meter} \\
    P5 & 0.40 & $-0.90$ & 1.55 & $+x$, \qty{0.60}{\meter} \\
    P6 & 1.00 & $-0.90$ & 1.55 & $+y$, \qty{0.80}{\meter} \\
    P7 & 1.00 & $-0.10$ & 1.55 & $-x$, \qty{0.60}{\meter} \\
    P8 & 0.40 & $-0.10$ & 1.55 & $-z$, \qty{0.25}{\meter} \\
    \bottomrule
  \end{tabularx}
\end{table}

\subsection{Time axis, state machine, and ISO legs}
\label{sec:sim-node}

Knot times are built on the \ocs{} relative axis
\texttt{observation.time}, never on the wall clock. A knot stamped in
wall time falls outside the receding horizon and the robot either freezes
or jumps to the last sample; that failure is cheap to see on the dummy
and expensive on hardware. At a trigger time $t_{0}$ the node walks the
landed corners, adding $\lVert \Delta\symbf{p}\rVert / v_{\mathrm{ref}}$
on every edge. The configured path speed is
$v_{\mathrm{ref}} = \qty{0.05}{\meter\per\second}$, half of the
\qty{0.10}{\meter\per\second} base limit in
\texttt{task\_experiment\_a.info}, so a dummy that cannot keep up is a
cost or constraint problem rather than a saturated plant. With
\texttt{corner\_dwell\_sec} $= \qty{2.0}{\second}$ the reference pauses
at each corner; the \qty{4.50}{\meter} perimeter then takes
\qty{90}{\second} of motion, and the seven dwells inside the measured
lap add \qty{14}{\second}, before lead-in and lead-out. A dwell of zero would
be the ISO path-accuracy mode (constant speed through corners). The
value used here is a pose-accuracy setting: the \ee{} is given time to
settle so that a later $\mathrm{AP}_{p}$ is not dominated by a corner
cut.

After \texttt{/start} the node does not immediately run the cycle. It
lands the box once, holds the reference at the lead-in point, and waits
until the measured \ee{} stays within
\qty{0.02}{\meter} of that point for \qty{0.3}{\second}
(\texttt{Approaching}). Only then does the running trajectory begin, so
the first measured edge does not contain the motion from standby. The
phase sequence is $\mathrm{Init}\to\mathrm{Hold}\to
\mathrm{Approaching}\to\mathrm{Running}\to\mathrm{Finished}$.
\texttt{auto\_start} is false: the box is drawn in RViz first, and a
service call starts the motion after a visual check.

ISO~9283 cl.~8.2 requires the programmed start and end to lie outside
the test path. A lead-in of \qty{0.15}{\meter} is placed on the backward
extension of P1--P2, and a lead-out of \qty{0.15}{\meter} on the tangent
of the closing edge. With the present Hamiltonian order that closing
edge is the downward $z$ leg, so the lead-out ends at
$z=\qty{1.15}{\meter}$. The dummy RViz configuration draws those legs
in orange for exactly this reason: a lead-out that would enter the
chassis is visible before anyone calls start. Returning to the standby
pose after the lead-out is switched off for this path
(\texttt{return\_to\_home: false}).

The trajectory is re-published at \qty{1}{\hz}. Re-sending the same
knots on the \ocs{} time axis is idempotent, but it restores the
reference if a reset handshake of the MRT bridge has overwritten it
with a two-sample hold (\cref{sec:hw-interfaces}). The dummy rarely
handshakes again in mid-run; the timer is still part of the node that
hardware will use.

Relative to \texttt{task\_hardware.info} (\cref{tab:mpc-config}), the
experiment file lowers the pre-scale diagonal of $R_{\mathrm{base}}$
from $8.5$ to $5.5$, replaces the uniform arm diagonal of $4.0$ by
per-joint values between $2.5$ and $4.0$, and raises the \ee{} position
weight from $10$ to $\mu_{\mathrm{p}}=16$ while lowering the orientation
weight from $5$ to $\mu_{\mathrm{o}}=3$. A cheaper base and a tighter
position term are what push the solver to reach the far corners by
driving the platform instead of stretching the arm towards its reach
limit. The file also lengthens the horizon from \qty{1.2}{\second} to
$T=\qty{1.3}{\second}$ and the solution time window from
\qty{1.0}{\second} to $T_{\mathrm{sol}}=\qty{1.2}{\second}$; solve and
evaluation rates stay at \qty{10}{\hz} and \qty{100}{\hz}. The dummy
therefore validates that operating point, not the everyday hardware
tune.

\subsection{Procedure and acceptance}
\label{sec:sim-acceptance}

A dummy session is started with
\texttt{experiment\_a\_dummy.launch.py}. After the first observation the
box appears in RViz, including the orange lead legs. The operator
confirms that every corner is in the cell and that the lead-out at
$z=\qty{1.15}{\meter}$ clears the chassis, then calls
\texttt{/ee\_waypoint\_target/start}. The dummy plant must reach the
entry point, traverse all eight measured edges, and finish on the
lead-out hold without the reference jumping to the last knot. The
metric topic must show \texttt{segment\_id} passing through $1$--$8$
and returning to a negative identifier. Those checks accept the
geometry, the \ocs{} time axis, the data pipeline and the visualisation.
They do not produce a tracking-accuracy number for the robot.

The same YAML and task file are the input of the mock-hardware and
physical runs in \cref{ch:hardware}. A planar four-corner variant
exists as a reduced path at constant $z$ and shares this node; it is
not used in the dummy checks reported here.

What the dummy cannot close is the list in
\cref{sec:sim-methodology}: policy lag, actuator dynamics, odometry
drift, and any ISO claim that needs an external \ee{} pose, $n\ge 3$
cycles, or a second pipeline. Those belong to the hardware chapter,
using the metrics of \cref{tab:sim-metrics} on the measured segments of
the path in \cref{fig:sim-box}.
